\chapter{General Context and Project Framework}

\section*{Introduction}

This chapter establishes the conceptual framework of the \caimp{} project.
We describe the technological environment of modern observability and the AIOps
domain, then formulate the problem precisely with supporting quantitative data
from industry and academic sources.

\section{Modern System Observability}

\subsection{The Three Pillars of Observability}

Observability refers to the ability to infer the internal state of a system
from its external outputs \cite{brendangregg2020}. In distributed computing,
it traditionally rests on three complementary pillars:

\begin{description}[leftmargin=2cm, style=nextline]
  \item[Metrics] Periodic numerical measurements of system state: CPU usage,
    memory occupation, request rate, latency. Their collection is lightweight
    and statistical analysis enables anomaly detection and trend forecasting.
  \item[Logs] Textual records of discrete events: application errors, HTTP
    requests, service starts. Rich in contextual information, they are
    indispensable for post-mortem diagnosis but difficult to analyse at scale
    without appropriate tooling.
  \item[Traces] Representations of a request's journey through the different
    components of a distributed system. They enable bottleneck localisation
    and dependency identification between services \cite{sigelman2010dapper}.
\end{description}

The \textsc{otlp} (\textit{OpenTelemetry Protocol}) \cite{otlp_spec},
standardised by the CNCF, unifies the collection of all three signal types
under a common format, facilitating interoperability between instrumentation
tools and analysis backends.

\subsection{From Reactive to Predictive Monitoring}

Early generations of monitoring tools, such as Nagios or Zabbix, adopted a
purely reactive approach: send an alert when a threshold is exceeded. This
philosophy had two major limitations. First, static thresholds generate a
significant proportion of false positives and false negatives on variable-load
systems. Second, the alert conveys no information about the \textit{cause}
of the problem.

The emergence of high-performance time-series databases like TimescaleDB
\cite{timescaledb2017} and statistical anomaly detection techniques
(Z-score, Isolation Forest \cite{liu2008isolation}, Holt-Winters
\cite{winters1960forecasting}) enabled a shift to a proactive paradigm:
detecting drift before it becomes an incident, and anticipating future
resource saturation.

\section{The AIOps Domain}

\subsection{Definition and Scope}

AIOps (\textit{Artificial Intelligence for IT Operations}) refers to the
application of machine learning and artificial intelligence techniques to
IT operations \cite{gartner_aiops2021}. Gartner, which coined the term in 2017,
defines it as ``a platform that combines Big Data, machine learning and advanced
analytical technologies to directly improve and automate IT operations''.

\subsection{Historical Evolution}

\begin{enumerate}
  \item \textbf{1990s--2000s}: Static threshold monitoring (Nagios, MRTG).
  \item \textbf{2000s--2010s}: Log and metric aggregation (Splunk, Graphite).
  \item \textbf{2010s--2018}: Distributed observability (Prometheus, Grafana, Jaeger).
  \item \textbf{2018--2023}: Emerging AIOps (Dynatrace Davis, Datadog Watchdog).
  \item \textbf{2023--present}: Generative AIOps (local LLMs, RAG). Explosion
    of open language models enabling natural-language explanation without
    cloud dependency.
\end{enumerate}

\section{Problem Statement}

\subsection{Alert Fatigue: A Documented Phenomenon}

Industry studies converge on identifying alert fatigue as one of the most
critical operational challenges. According to the PagerDuty 2023 report
\cite{pagerduty2023}:

\begin{itemize}
  \item 73\% of engineers report regularly ignoring alerts due to excessive volume;
  \item 40\% of alerts are false positives or duplicates;
  \item Average MTTR is 4 hours 20 minutes for organisations without AIOps
    tooling, versus 42 minutes for those that have adopted it.
\end{itemize}

\subsection{The Gap Between Detection and Understanding}

The standard response to a ``CPU at 98\%'' alert follows an implicit
investigation protocol that any operations engineer will recognise:
check consumer processes, consult application logs, examine recent deployment
history, check slow database queries, consult network metrics, and cross-reference
with system events. This sequence, repeated at 2am during on-call duty, is a
source of errors and stress, requiring expertise that small teams do not
systematically possess.

\subsection{Problem Formulation}

\begin{infobox}[Problem Statement]
\textit{How can we design a self-hosted infrastructure monitoring platform,
accessible to teams without DevOps expertise, capable not only of detecting
anomalies but of explaining their cause in natural language and proposing
concrete corrective actions, all without dependency on third-party cloud services?}
\end{infobox}

\section{Functional Scope}

\textbf{In scope:} CPU/memory/disk metric collection via OTLP agent; three-algorithm
anomaly detection; change correlation root cause analysis; 24-hour Holt-Winters
forecasting; answers-first dashboard; alert deduplication; PDF reports;
optional Splunk integration; Docker Compose deployment; mTLS and JWT security.

\textbf{Out of scope:} Distributed tracing; automated remediation execution;
Windows agent support; mobile interface; federated authentication (SAML, OIDC).

\section*{Conclusion}

This chapter defined the conceptual framework of the \caimp{} project, anchoring
the problem in the real challenges of modern monitoring. The analysis of alert
fatigue and the gap between detection and understanding fully justifies the AIOps
approach adopted. The following chapter surveys existing solutions and positions
\caimp{} within this ecosystem.
